\section{Circuit Complexity --- Meyer Machine Parse}
%%% generated by the blueprint pipeline from TCSlib.Complexity.CircuitComplexity.MeyerMachineParse
\subsection*{Circuit Complexity --- Meyer Machine Parse: Introduction}

Declarations of the Lean module \texttt{TCSlib.Complexity.CircuitComplexity.MeyerMachineParse}.

\subsection*{Circuit Complexity --- Meyer Machine Parse: Results}

\begin{definition}[Boolean separator code]
\lean{Complexity.Meyer.fcode}
\label{Complexity.Meyer.fcode}
\leanok
For a finite Boolean string $w$, define its separator code by replacing each bit $b$ by
the two-bit string $(b,\mathrm{true})$ and concatenating these pairs in order.
\end{definition}

\begin{theorem}[Length of the encoded word]
\lean{Complexity.Meyer.length_fcode}
\label{Complexity.Meyer.length_fcode}
\leanok
\uses{Complexity.Meyer.fcode}
\difficulty{4}
For every Boolean word $w$, its encoding, formed by replacing each bit $b$ with the pair
$(b,1)$, has length $2|w|$.
\end{theorem}

\begin{definition}[Reading a Delimited Boolean Field]
\lean{Complexity.Meyer.readField}
\label{Complexity.Meyer.readField}
\leanok
For a Boolean list, read a field by consuming pairs whose second bit is true: each such
pair contributes its first bit to the field. A pair whose second bit is false terminates
the field, returning the accumulated bits and the remaining list after that pair; if the
list ends before a terminator is found, the result is undefined.
\end{definition}

\begin{theorem}[Reading a Coded Field]
\lean{Complexity.Meyer.readField_fcode}
\label{Complexity.Meyer.readField_fcode}
\leanok
\uses{Complexity.Meyer.fcode, Complexity.Meyer.readField}
\difficulty{4}
For any bit strings $w$ and $r$ and any bit $e$, reading a field from the field code of
$w$ followed by $e$, then the delimiter $0$, and then $r$ returns exactly $(w,r)$.
\end{theorem}

\begin{theorem}[Structure of a Successfully Read Field]
\lean{Complexity.Meyer.readField_eq_some}
\label{Complexity.Meyer.readField_eq_some}
\leanok
\uses{Complexity.Meyer.fcode, Complexity.Meyer.readField}
\difficulty{5}
For bit strings $s$, $w$, and $r$, if reading a field from $s$ returns the pair $(w,r)$,
then there is a bit $e$ such that $s$ is the field code of $w$, followed by $e$, then
the delimiter $0$, and then $r$.
\end{theorem}

\begin{definition}[Encoding a Query String]
\lean{Complexity.Meyer.query}
\label{Complexity.Meyer.query}
\leanok
\uses{Complexity.Meyer.fcode}
Let $n$ be a natural number, let $b_0,\ldots,b_{n-1}$ be Boolean bits, and let
$t,a,x$ be Boolean lists. The query string is the bit list
$(b_0,\ldots,b_{n-1})$ followed by the code of $t$, two false bits, the code of
$a$, two false bits, and $x$.
\end{definition}

\begin{theorem}[Length of a query word]
\lean{Complexity.Meyer.length_query}
\label{Complexity.Meyer.length_query}
\leanok
\uses{Complexity.Meyer.query}
\difficulty{3}
Let $n_b$ be a natural number, let $b$ be a Boolean string of length $n_b$, and let
$t,a,x$ be Boolean words. The query formed by concatenating $b$, the encodings of $t$
and $a$ separated by two zero bits, and then $x$ after two further zero bits has length
$n_b+2|t|+2|a|+4+|x|$.
\end{theorem}

\begin{definition}[Parsing a Query String]
\lean{Complexity.Meyer.parseQ}
\label{Complexity.Meyer.parseQ}
\leanok
\uses{Complexity.Meyer.readField}
Let $n$ be a natural number and $y$ a Boolean list. If $y$ has at least $n$ entries,
take its first $n$ bits as a function on the indices $0,\ldots,n-1$, and parse
the remainder as two consecutive delimited fields using the Boolean-field reader. When
both fields parse successfully, return the first $n$ bits and the two parsed fields
together with the unconsumed suffix; otherwise return no result.
\end{definition}

\begin{theorem}[Parsing an Encoded Query]
\lean{Complexity.Meyer.parseQ_query}
\label{Complexity.Meyer.parseQ_query}
\leanok
\uses{Complexity.Meyer.fcode, Complexity.Meyer.length_query, Complexity.Meyer.parseQ, Complexity.Meyer.query, Complexity.Meyer.readField_fcode}
\difficulty{5}
For any natural number $n$, bit vector
$b\colon\{0,\ldots,n-1\}\to\{\mathrm{false},\mathrm{true}\}$, and Boolean words $t,a,x$,
encode them as the concatenation of $b$ with the delimiter encoding of $t$, then two
false bits, then the delimiter encoding of $a$, then two false bits, and finally $x$.
Parsing this word with bit-vector length $n$ returns exactly $(b,t,a,x)$.
\end{theorem}

\begin{definition}[Clamped Input Position]
\lean{Complexity.Meyer.fpos}
\label{Complexity.Meyer.fpos}
\leanok
Let $M$ be a finite Turing machine over the Boolean alphabet, let $y$ be a Boolean list,
and let $p$ be a natural number. The position is the index $\min(p, |y|+1)$ in the
range from $0$ to $|y|+1$.
\end{definition}

\begin{theorem}[Rightward movement of a virtual input position]
\lean{Complexity.Meyer.moveInputPos_fpos}
\label{Complexity.Meyer.moveInputPos_fpos}
\leanok
\uses{Complexity.Meyer.fpos, Turing.moveInputPos, Turing.moveInputPos_pos_of_ne_right}
\difficulty{4}
Let $y$ be a Boolean word and let $p$ be a natural number with $p\le |y|$. If the
virtual input-head position is the position corresponding to $p$, then moving it one
step to the right gives the position corresponding to $p+1$.
\end{theorem}

\begin{theorem}[Input symbol at a clamped position]
\lean{Complexity.Meyer.inputSymbol_fpos}
\label{Complexity.Meyer.inputSymbol_fpos}
\leanok
\uses{Complexity.Meyer.fpos, Turing.Cfg, Turing.Cfg.inputSymbol, Turing.FinTM.inputSymbol_at}
\difficulty{3}
Let $c$ be a configuration with Boolean input $y$, and let $p\in\mathbb{N}$ satisfy
$1\le p\le |y|+1$. If the input head of $c$ is at position $p$, clamped to the valid
range from $0$ through $|y|+1$, then the symbol it reads is the possibly undefined list
entry $y[p-1]$.
\end{theorem}

\begin{definition}[Auxiliary Machine Configuration]
\lean{Complexity.Meyer.acfg}
\label{Complexity.Meyer.acfg}
\leanok
\uses{Complexity.Meyer.St, Complexity.Meyer.tabTM, Complexity.Meyer.tcfg, Turing.Cfg}
Let $M$ be a finite Turing machine over the Boolean alphabet, let $y$ be its input, let
$s$ be a state of the tabulating machine built from $M$, and let $p$ be an input-head
position among the $|y|+2$ positions flanking the input. Choose three integer-indexed
Boolean tapes $\tau_0,\tau_1,\tau_2$, their head positions $z_0,z_1,z_2\in\mathbb Z$,
and an output list. The resulting configuration has state $s$, input-head position $p$,
blank contents and zero head positions on the original work tapes, and the chosen
contents, head positions, and output on the three auxiliary tapes and output tape.
\end{definition}

\begin{theorem}[Step on an auxiliary configuration]
\lean{Complexity.Meyer.acfg_step}
\label{Complexity.Meyer.acfg_step}
\leanok
\uses{Complexity.Meyer.St, Complexity.Meyer.acfg, Complexity.Meyer.tabTM, Complexity.Meyer.tabTM_step, Complexity.Meyer.tabTr, Turing.Action.apply, Turing.Cfg.inputSymbol, Turing.MultiTapeTM.step}
\difficulty{4}
Let $M$ be a finite-state Turing machine with Boolean alphabet and let $y$ be a finite
Boolean input. For any simulated state $q$, input-head position $p$ among the $|y|+2$
positions, three auxiliary tape contents and head positions, and output word, the
tabulating machine's step from the auxiliary configuration with state $q$ equals the
result of applying its tabulated transition for $q$, using the scanned input symbol and
the symbols read from the three auxiliary tapes (with all simulated work tapes blank).
\end{theorem}

\begin{theorem}[Effect with Idle Original Tapes]
\lean{Complexity.Meyer.tact_idle_apply}
\label{Complexity.Meyer.tact_idle_apply}
\leanok
\uses{Complexity.Meyer.St, Complexity.Meyer.acfg, Complexity.Meyer.idle, Complexity.Meyer.tact, Complexity.Meyer.tact_apply, Complexity.Meyer.wr, Turing.Action.apply, Turing.moveInputPos}
\difficulty{4}
For a Boolean finite Turing machine $M$, any Boolean input word, input-head movement,
instructions for the three auxiliary tapes, optional output bit and successor state, and
configuration data for the input head, auxiliary tapes, and output, applying the action
that leaves every original work tape unchanged yields the configuration with the
prescribed successor state, the input head moved by the clamped movement, each auxiliary
tape updated at its head when instructed and each auxiliary head moved as prescribed,
and the optional output bit appended.
\end{theorem}

\begin{theorem}[Idle action entry]
\lean{Complexity.Meyer.idle_apply}
\label{Complexity.Meyer.idle_apply}
\leanok
\uses{Complexity.Meyer.idle}
\difficulty{1}
For every natural number $n$ and index $j<n$, the $j$th entry of the idle action table
is the pair consisting of no write instruction and zero movement.
\end{theorem}

\begin{theorem}[Value at the Updated Coordinate]
\lean{Complexity.Meyer.one_0_0}
\label{Complexity.Meyer.one_0_0}
\leanok
\uses{Complexity.Meyer.one}
\difficulty{1}
Let $M$ be a finite Turing machine with Boolean symbols, and let $a$ be a pair
consisting of an optional optional Boolean and a sign. For the function on three
coordinates that equals $a$ at coordinate $0$ and is otherwise idle, its value at
coordinate $0$ is $a$.
\end{theorem}

\begin{theorem}[Idle Value Away from the Updated Coordinate]
\lean{Complexity.Meyer.one_0_1}
\label{Complexity.Meyer.one_0_1}
\leanok
\uses{Complexity.Meyer.one}
\difficulty{1}
Let $M$ be a finite Turing machine with Boolean symbols, and let $a$ be a pair
consisting of an optional optional Boolean and a sign. For the function on three
coordinates that equals $a$ at coordinate $0$ and is otherwise idle, its value at
coordinate $1$ is $(	ext\{none\},0)$.
\end{theorem}

\begin{theorem}[Idle Value Away from the Updated Coordinate]
\lean{Complexity.Meyer.one_0_2}
\label{Complexity.Meyer.one_0_2}
\leanok
\uses{Complexity.Meyer.one}
\difficulty{1}
Let $M$ be a finite Turing machine with Boolean symbols, and let $a$ be a pair
consisting of an optional optional Boolean and a sign. For the function on three
coordinates that equals $a$ at coordinate $0$ and is otherwise idle, its value at
coordinate $2$ is $(	ext\{none\},0)$.
\end{theorem}

\begin{theorem}[Idle Value Away from the Updated Coordinate]
\lean{Complexity.Meyer.one_1_0}
\label{Complexity.Meyer.one_1_0}
\leanok
\uses{Complexity.Meyer.one}
\difficulty{1}
Let $M$ be a finite Turing machine with Boolean symbols, and let $a$ be a pair
consisting of an optional optional Boolean and a sign. For the function on three
coordinates that equals $a$ at coordinate $1$ and is otherwise idle, its value at
coordinate $0$ is $(	ext\{none\},0)$.
\end{theorem}

\begin{theorem}[Value at the Updated Coordinate]
\lean{Complexity.Meyer.one_1_1}
\label{Complexity.Meyer.one_1_1}
\leanok
\uses{Complexity.Meyer.one}
\difficulty{1}
Let $M$ be a finite Turing machine with Boolean symbols, and let $a$ be a pair
consisting of an optional optional Boolean and a sign. For the function on three
coordinates that equals $a$ at coordinate $1$ and is otherwise idle, its value at
coordinate $1$ is $a$.
\end{theorem}

\begin{theorem}[Idle Value Away from the Updated Coordinate]
\lean{Complexity.Meyer.one_1_2}
\label{Complexity.Meyer.one_1_2}
\leanok
\uses{Complexity.Meyer.one}
\difficulty{1}
Let $M$ be a finite Turing machine with Boolean symbols, and let $a$ be a pair
consisting of an optional optional Boolean and a sign. For the function on three
coordinates that equals $a$ at coordinate $1$ and is otherwise idle, its value at
coordinate $2$ is $(	ext\{none\},0)$.
\end{theorem}

\begin{theorem}[Idle Value Away from the Updated Coordinate]
\lean{Complexity.Meyer.one_2_0}
\label{Complexity.Meyer.one_2_0}
\leanok
\uses{Complexity.Meyer.one}
\difficulty{1}
Let $M$ be a finite Turing machine with Boolean symbols, and let $a$ be a pair
consisting of an optional optional Boolean and a sign. For the function on three
coordinates that equals $a$ at coordinate $2$ and is otherwise idle, its value at
coordinate $0$ is $(	ext\{none\},0)$.
\end{theorem}

\begin{theorem}[Idle Value Away from the Updated Coordinate]
\lean{Complexity.Meyer.one_2_1}
\label{Complexity.Meyer.one_2_1}
\leanok
\uses{Complexity.Meyer.one}
\difficulty{1}
Let $M$ be a finite Turing machine with Boolean symbols, and let $a$ be a pair
consisting of an optional optional Boolean and a sign. For the function on three
coordinates that equals $a$ at coordinate $2$ and is otherwise idle, its value at
coordinate $1$ is $(	ext\{none\},0)$.
\end{theorem}

\begin{theorem}[Value at the Updated Coordinate]
\lean{Complexity.Meyer.one_2_2}
\label{Complexity.Meyer.one_2_2}
\leanok
\uses{Complexity.Meyer.one}
\difficulty{1}
For a three-coordinate function whose default value is idle and whose value at
coordinate $2$ is replaced by $a$, the value at coordinate $2$ is $a$.
\end{theorem}

\begin{theorem}[First updated coordinate]
\lean{Complexity.Meyer.two_01_0}
\label{Complexity.Meyer.two_01_0}
\leanok
\uses{Complexity.Meyer.two}
\difficulty{1}
For arbitrary values $a,a' \in \mathrm{Option}(\mathrm{Option}(\{false,true\}))\times \mathrm{SignType}$, the
function on the three coordinates $0,1,2$ obtained by first setting coordinate $0$ to
$a$ and then coordinate $1$ to $a'$ takes the value $a$ at coordinate $0$.
\end{theorem}

\begin{theorem}[Second updated coordinate]
\lean{Complexity.Meyer.two_01_1}
\label{Complexity.Meyer.two_01_1}
\leanok
\uses{Complexity.Meyer.two}
\difficulty{1}
For arbitrary values $a,a' \in \mathrm{Option}(\mathrm{Option}(\{false,true\}))\times \mathrm{SignType}$, the
function on the three coordinates $0,1,2$ obtained by first setting coordinate $0$ to
$a$ and then coordinate $1$ to $a'$ takes the value $a'$ at coordinate $1$.
\end{theorem}

\begin{theorem}[Unmodified third coordinate]
\lean{Complexity.Meyer.two_01_2}
\label{Complexity.Meyer.two_01_2}
\leanok
\uses{Complexity.Meyer.two}
\difficulty{1}
For arbitrary values $a,a' \in \mathrm{Option}(\mathrm{Option}(\{false,true\}))\times \mathrm{SignType}$, the
function on the three coordinates $0,1,2$ obtained by first setting coordinate $0$ to
$a$ and then coordinate $1$ to $a'$ takes the default value $(\{none\},0)$ at coordinate
$2$.
\end{theorem}

\begin{theorem}[Unmodified first coordinate]
\lean{Complexity.Meyer.two_12_0}
\label{Complexity.Meyer.two_12_0}
\leanok
\uses{Complexity.Meyer.two}
\difficulty{1}
For arbitrary values $a,a' \in \mathrm{Option}(\mathrm{Option}(\{false,true\}))\times \mathrm{SignType}$, the
function on the three coordinates $0,1,2$ obtained by first setting coordinate $1$ to
$a$ and then coordinate $2$ to $a'$ takes the default value $(\{none\},0)$ at coordinate
$0$.
\end{theorem}

\begin{theorem}[Value at the First Updated Position]
\lean{Complexity.Meyer.two_12_1}
\label{Complexity.Meyer.two_12_1}
\leanok
\uses{Complexity.Meyer.two}
\difficulty{1}
Let $a$ and $a'$ be values in $\mathrm{Option}(\mathrm{Option}(\{0,1\}))\times\{
-1,0,1\}$. On the three positions $1,2,3$, start with the constant default value, set
position $1$ to $a$, and then set position $2$ to $a'$. The resulting function has value
$a$ at position $1$.
\end{theorem}

\begin{theorem}[Value at the Second Updated Position]
\lean{Complexity.Meyer.two_12_2}
\label{Complexity.Meyer.two_12_2}
\leanok
\uses{Complexity.Meyer.two}
\difficulty{1}
Let $a$ and $a'$ be values in $\mathrm{Option}(\mathrm{Option}(\{0,1\}))\times\{
-1,0,1\}$. On the three positions $1,2,3$, start with the constant default value, set
position $1$ to $a$, and then set position $2$ to $a'$. The resulting function has value
$a'$ at position $2$.
\end{theorem}

\begin{theorem}[First coordinate after updates at zero and two]
\lean{Complexity.Meyer.two_02_0}
\label{Complexity.Meyer.two_02_0}
\leanok
\uses{Complexity.Meyer.two}
\difficulty{1}
For arbitrary values $a,a' \in \mathrm{Option}(\mathrm{Option}(\{false,true\}))\times \mathrm{SignType}$, the
function on the three coordinates $0,1,2$ obtained by first setting coordinate $0$ to
$a$ and then coordinate $2$ to $a'$ takes the value $a$ at coordinate $0$.
\end{theorem}

\begin{theorem}[Unmodified middle coordinate]
\lean{Complexity.Meyer.two_02_1}
\label{Complexity.Meyer.two_02_1}
\leanok
\uses{Complexity.Meyer.two}
\difficulty{1}
For arbitrary values $a,a' \in \mathrm{Option}(\mathrm{Option}(\{false,true\}))\times \mathrm{SignType}$, the
function on the three coordinates $0,1,2$ obtained by first setting coordinate $0$ to
$a$ and then coordinate $2$ to $a'$ takes the default value $(\{none\},0)$ at coordinate
$1$.
\end{theorem}

\begin{theorem}[Second updated coordinate]
\lean{Complexity.Meyer.two_02_2}
\label{Complexity.Meyer.two_02_2}
\leanok
\uses{Complexity.Meyer.two}
\difficulty{1}
For arbitrary values $a,a' \in \mathrm{Option}(\mathrm{Option}(\{false,true\}))\times \mathrm{SignType}$, the
function on the three coordinates $0,1,2$ obtained by first setting coordinate $0$ to
$a$ and then coordinate $2$ to $a'$ takes the value $a'$ at coordinate $2$.
\end{theorem}

\begin{theorem}[No-Write Tape Update]
\lean{Complexity.Meyer.wr_none}
\label{Complexity.Meyer.wr_none}
\leanok
\uses{Complexity.Meyer.wr}
\difficulty{1}
For any tape contents $\tau:\mathbb Z\to\text{Option}(\text{Bool})$ and cell position
$z\in\mathbb Z$, applying the write operation with no write prescribed leaves $\tau$
unchanged.
\end{theorem}

\begin{theorem}[Prescribed Tape Write]
\lean{Complexity.Meyer.wr_some}
\label{Complexity.Meyer.wr_some}
\leanok
\uses{Complexity.Meyer.wr}
\difficulty{1}
For any optional Boolean symbol $s$, tape contents $\tau:\mathbb
Z\to\text{Option}(\text{Bool})$, and cell position $z\in\mathbb Z$, applying the write
operation with $s$ prescribed gives the tape contents obtained by replacing the value at
$z$ by $s$ and leaving every other cell unchanged.
\end{theorem}

\begin{definition}[Blank Boolean Tape]
\lean{Complexity.Meyer.bl}
\label{Complexity.Meyer.bl}
\leanok
For a finite Turing machine $M$ over the Boolean alphabet, the blank tape is the
integer-indexed tape that contains no symbol at any position.
\end{definition}

\begin{theorem}[Initial configuration of the tabulating machine]
\lean{Complexity.Meyer.init_eq}
\label{Complexity.Meyer.init_eq}
\leanok
\uses{Complexity.Meyer.acfg, Complexity.Meyer.bl, Complexity.Meyer.fpos, Complexity.Meyer.nbitsOf_pos, Complexity.Meyer.tabTM, Complexity.Meyer.tcfg, Turing.MultiTapeTM.initCfg}
\difficulty{4}
Let $M$ be a finite Turing machine over the Boolean alphabet, and let $y$ be a finite
Boolean input. The initial configuration of the machine constructed from $M$ has its
simulated state register set to the initial state of $M$ with an empty output register,
its input head at position $1$ clamped to the range of positions for $y$, all work-tape
cells blank, all work-tape heads at the origin, and an empty output tape.
\end{theorem}

\begin{definition}[Selected Bit Prefix]
\lean{Complexity.Meyer.selPart}
\label{Complexity.Meyer.selPart}
\leanok
\uses{Complexity.Meyer.nbits}
Let $M$ be a finite Turing machine over the Boolean alphabet, let $y$ be a Boolean list,
and let $m$ be a natural number. Define a Boolean string of length equal to the number
of selector types for $M$ by taking, at index $j$, the $j$th bit of $y$ when $j<m$,
using false if that index is beyond the end of $y$, and taking false when $j\ge m$.
\end{definition}

\begin{theorem}[Successive Selector Prefixes]
\lean{Complexity.Meyer.selPart_succ}
\label{Complexity.Meyer.selPart_succ}
\leanok
\uses{Complexity.Meyer.nbits, Complexity.Meyer.selPart}
\difficulty{4}
Let $M$ be a finite Boolean Turing machine, let $y$ be a Boolean word, and let $m$ be a
natural number with $m<\mathrm{nbits}(M)$ and $m<|y|$. Updating the $m$-bit selector
prefix of $y$ at coordinate $m$ to the bit $y_m$ gives the $(m+1)$-bit selector prefix
of $y$.
\end{theorem}

\begin{theorem}[Selector Phase of the Tabulation]
\lean{Complexity.Meyer.sel_run}
\label{Complexity.Meyer.sel_run}
\leanok
\uses{Complexity.Meyer.acfg, Complexity.Meyer.acfg_step, Complexity.Meyer.bl, Complexity.Meyer.fpos, Complexity.Meyer.idle_apply, Complexity.Meyer.init_eq, Complexity.Meyer.inputSymbol_fpos, Complexity.Meyer.moveInputPos_fpos, Complexity.Meyer.nbits, Complexity.Meyer.selPart, Complexity.Meyer.selPart_succ, Complexity.Meyer.tabTM, Complexity.Meyer.tabTr, Complexity.Meyer.tact_idle_apply, Complexity.Meyer.wr_none, Turing.MultiTapeTM.initCfg, Turing.MultiTapeTM.runFrom, Turing.MultiTapeTM.runFrom_succ_eq_step', Turing.MultiTapeTM.runFrom_zero}
\difficulty{5}
Let $M$ be a finite Turing machine with Boolean input alphabet, and write $n$ for the
number of selector states in its tabulation. For every Boolean string $y$, every $m<n$,
and every $m\le |y|$, after $m$ steps the tabulated machine is in selector state $m$,
with the first $m$ selector bits equal to the corresponding bits of $y$ and all
remaining bits false. Its input head is at position $\min(m+1,|y|+1)$ in the input tape
convention with boundary blanks; all work tapes and the three auxiliary tapes are blank,
their heads are at the origin, and the output is empty.
\end{theorem}

\begin{theorem}[Completion of the Selector Scan]
\lean{Complexity.Meyer.sel_done}
\label{Complexity.Meyer.sel_done}
\leanok
\uses{Complexity.Meyer.acfg, Complexity.Meyer.acfg_step, Complexity.Meyer.bl, Complexity.Meyer.fpos, Complexity.Meyer.idle_apply, Complexity.Meyer.inputSymbol_fpos, Complexity.Meyer.moveInputPos_fpos, Complexity.Meyer.nbits, Complexity.Meyer.nbitsOf_pos, Complexity.Meyer.selPart, Complexity.Meyer.selPart_succ, Complexity.Meyer.sel_run, Complexity.Meyer.tabTM, Complexity.Meyer.tabTr, Complexity.Meyer.tact_idle_apply, Complexity.Meyer.wr_none, Turing.MultiTapeTM.initCfg, Turing.MultiTapeTM.runFrom, Turing.MultiTapeTM.runFrom_succ_eq_step'}
\difficulty{5}
Let $M$ be a finite Boolean Turing machine and $y$ a Boolean input word whose length is
at least the number of selector bits for $M$. After that many steps of $M$'s tabulating
machine from its initial configuration on $y$, the machine is in state $\mathrm{fA}$
with false flag and selector bits given by the first $\mathrm{nbits}(M)$ bits of $y$.
Its input head is at position $\min(\mathrm{nbits}(M)+1,|y|+1)$; all three auxiliary
tapes are blank with their heads at zero, all simulated work tapes are blank with heads
at zero, and the output is empty.
\end{theorem}

\begin{theorem}[Rejection on a Short Selector String]
\lean{Complexity.Meyer.sel_reject}
\label{Complexity.Meyer.sel_reject}
\leanok
\uses{Complexity.Meyer.acfg, Complexity.Meyer.acfg_step, Complexity.Meyer.inputSymbol_fpos, Complexity.Meyer.nbits, Complexity.Meyer.reject, Complexity.Meyer.sel_run, Complexity.Meyer.tabTM, Complexity.Meyer.tabTr, Complexity.Meyer.tact_idle_apply, Complexity.Meyer.tcfg, Turing.MultiTapeTM.initCfg, Turing.MultiTapeTM.runFrom, Turing.MultiTapeTM.runFrom_succ_eq_step'}
\difficulty{4}
Let $M$ be a finite Turing machine with Boolean input alphabet, and let $n$ be the
number of selector symbols in its finite tabulation. For every Boolean string $y$ with
$|y|<n$, the tabulated machine, started on $y$, is halted after $|y|+1$ steps and has
output consisting of the single bit $0$.
\end{theorem}

\begin{definition}[Configuration from Three Buffers]
\lean{Complexity.Meyer.pcfg}
\label{Complexity.Meyer.pcfg}
\leanok
\uses{Complexity.Meyer.St, Complexity.Meyer.acfg, Complexity.Meyer.tabTM, Turing.Cfg, Turing.FinTM.bufferTape}
Let $M$ be a finite Turing machine over the Boolean alphabet, let $y$ be an input list,
let $s$ be a state of the tabulating machine built from $M$, and let $p$ be an
input-head position among the $|y|+2$ positions. Given three Boolean lists
$w_X,w_C,w_O$, form the configuration with state $s$, input position $p$, blank original
work tapes at position zero, and auxiliary tapes encoding these three lists with their
heads at positions $|w_X|,|w_C|,|w_O|$, respectively; the output is empty.
\end{definition}

\begin{theorem}[Step of the encoded simulator]
\lean{Complexity.Meyer.pcfg_step}
\label{Complexity.Meyer.pcfg_step}
\leanok
\uses{Complexity.Meyer.St, Complexity.Meyer.acfg_step, Complexity.Meyer.fpos, Complexity.Meyer.inputSymbol_fpos, Complexity.Meyer.pcfg, Complexity.Meyer.tabTM, Complexity.Meyer.tabTr, Turing.Action.apply, Turing.FinTM.bufferTape, Turing.MultiTapeTM.step}
\difficulty{3}
Let $M$ be a finite Turing machine over the Boolean alphabet, let $y$ be a Boolean input
word, let $st$ be a state of the simulator constructed from $M$, and let $w_X,w_C,w_O$
be Boolean words. For a natural number $p$ with $1\le p\le |y|+1$, form the simulator
configuration with state $st$, input-head position clipped to the valid range, auxiliary
tapes containing these three words, and all original work tapes blank. Its next
configuration is exactly the result of applying the simulator transition for state $st$,
input symbol $y[p-1]?$ (with no symbol when that index is out of range), and the symbols
scanned at the right ends of the three auxiliary words, with every original work-tape
symbol blank.
\end{theorem}

\begin{theorem}[First blank cell after a stored word]
\lean{Complexity.Meyer.bufferTape_length}
\label{Complexity.Meyer.bufferTape_length}
\leanok
\uses{Turing.FinTM.bufferTape, Turing.FinTM.bufferTape_nat}
\difficulty{3}
For every Boolean word $w$, the tape storing $w$ starting at cell zero is blank at cell
$|w|$.
\end{theorem}

\begin{theorem}[Appending a symbol to a stored word]
\lean{Complexity.Meyer.bufferTape_write}
\label{Complexity.Meyer.bufferTape_write}
\leanok
\uses{Turing.FinTM.bufferTape, Turing.FinTM.bufferTape_append}
\difficulty{1}
Let $w$ be a Boolean word and let $b$ be a Boolean value. Updating the tape that stores
$w$ at cell $|w|$ to contain $b$ yields the tape that stores the concatenation of $w$
with the one-symbol word $[b]$.
\end{theorem}

\begin{theorem}[Rejection on a Blank Input Symbol]
\lean{Complexity.Meyer.pcfg_reject}
\label{Complexity.Meyer.pcfg_reject}
\leanok
\uses{Complexity.Meyer.St, Complexity.Meyer.acfg, Complexity.Meyer.fpos, Complexity.Meyer.pcfg, Complexity.Meyer.pcfg_step, Complexity.Meyer.reject, Complexity.Meyer.tabTM, Complexity.Meyer.tabTr, Complexity.Meyer.tact_idle_apply, Complexity.Meyer.tcfg, Turing.MultiTapeTM.step}
\difficulty{3}
Let $M$ be a finite Turing machine over Booleans and let $y$ be a Boolean input word.
Suppose that, for every work-tape contents function, the constructed transition rule in
state $st$ on a blank input symbol is the rejecting action. Let $p$ satisfy $1\le p\le
|y|+1$, and suppose the symbol at zero-based position $p-1$ is absent. For any three
auxiliary words, one step from the corresponding configuration has no successor state
and has output exactly the one-bit word $[\mathrm{false}]$.
\end{theorem}

\begin{theorem}[Transition from the First Field State]
\lean{Complexity.Meyer.pcfg_fA}
\label{Complexity.Meyer.pcfg_fA}
\leanok
\uses{Complexity.Meyer.fpos, Complexity.Meyer.moveInputPos_fpos, Complexity.Meyer.nbits, Complexity.Meyer.pcfg, Complexity.Meyer.pcfg_step, Complexity.Meyer.tabTM, Complexity.Meyer.tabTr, Complexity.Meyer.tact_idle_apply, Turing.MultiTapeTM.step}
\difficulty{3}
Let $M$ be a finite Turing machine over Booleans, let $y$ be a Boolean input word, and
let $p$ be a natural number with $1\le p\le |y|$. Suppose the symbol of $y$ at
zero-based position $p-1$ is $b$. For any flag $f$, bit vector, and three Boolean words
$w_X,w_C,w_O$, one step of the constructed machine from the configuration in state $fA$
with input head at position $p$ and those words on its auxiliary tapes yields the
configuration in state $fB$ with flag $f$, bit $b$, and the same bit vector, words, and
output; the input head advances to position $p+1$.
\end{theorem}

\begin{theorem}[True Branch from the Second Field State]
\lean{Complexity.Meyer.pcfg_fB_true}
\label{Complexity.Meyer.pcfg_fB_true}
\leanok
\uses{Complexity.Meyer.bufferTape_write, Complexity.Meyer.fpos, Complexity.Meyer.moveInputPos_fpos, Complexity.Meyer.nbits, Complexity.Meyer.pcfg, Complexity.Meyer.pcfg_step, Complexity.Meyer.tabTM, Complexity.Meyer.tabTr, Complexity.Meyer.tact_idle_apply, Turing.MultiTapeTM.step}
\difficulty{4}
Let $M$ be a finite Turing machine over Booleans, let $y$ be a Boolean input word, and
let $p$ satisfy $1\le p\le |y|$. Suppose the symbol of $y$ at zero-based position $p-1$
is true. For any Boolean flags $f,b$, bit vector, and auxiliary words $w_X,w_C,w_O$, one
step of the constructed machine from state $fB$ produces state $fA$ with flag $f$ and
the same bit vector, while advancing the input head to position $p+1$. If $f$ is true,
it appends $b$ to $w_O$ and leaves $w_C$ unchanged; if $f$ is false, it appends $b$ to
$w_C$ and leaves $w_O$ unchanged. The word $w_X$ is unchanged in either case.
\end{theorem}

\begin{theorem}[False Branch from the Second Field State]
\lean{Complexity.Meyer.pcfg_fB_false}
\label{Complexity.Meyer.pcfg_fB_false}
\leanok
\uses{Complexity.Meyer.fpos, Complexity.Meyer.moveInputPos_fpos, Complexity.Meyer.nbits, Complexity.Meyer.pcfg, Complexity.Meyer.pcfg_step, Complexity.Meyer.tabTM, Complexity.Meyer.tabTr, Complexity.Meyer.tact_idle_apply, Turing.MultiTapeTM.step}
\difficulty{3}
Let $M$ be a finite Turing machine over Booleans, let $y$ be a Boolean input word, and
let $p$ satisfy $1\le p\le |y|$. Suppose the symbol of $y$ at zero-based position $p-1$
is false. For any Boolean flags $f,b$, bit vector, and auxiliary words $w_X,w_C,w_O$,
one step of the constructed machine from state $fB$ with input head at position $p$
produces state $cpy$ with the same bit vector and auxiliary words if $f$ is true, and
state $fA$ with flag true and the same data if $f$ is false; the input head advances to
position $p+1$.
\end{theorem}
